\documentclass[11pt,a4paper]{article}
\usepackage{fontspec,xeCJK}
\setCJKmainfont{FandolSong-Regular}[Extension=.otf,BoldFont=FandolSong-Bold,ItalicFont=FandolKai-Regular]
\setCJKsansfont{FandolHei-Regular}[Extension=.otf,BoldFont=FandolHei-Bold]
\setCJKmonofont{FandolFang-Regular}[Extension=.otf]
\renewcommand{\abstractname}{摘要}
\renewcommand{\refname}{参考文献}
\linespread{1.12}
\usepackage{microtype}
\usepackage[margin=.85in]{geometry}
\usepackage{amsmath,amssymb,amsthm,mathtools,aliascnt}
\usepackage{booktabs,tabularx,array,enumitem,needspace,float}
\usepackage[section]{placeins}
\usepackage{xcolor,tikz,mdframed}
\usetikzlibrary{arrows.meta,positioning,calc,decorations.pathreplacing}
\usepackage[hidelinks]{hyperref}
\usepackage[nameinlink,noabbrev]{cleveref}
\usepackage{fancyhdr}
\hypersetup{pdftitle={Sync HotStuff RR v4.3},pdfsubject={区域终局性与 L1 授权轮换}}
\setlength{\headheight}{14pt}
\pagestyle{fancy}\fancyhf{}
\lhead{\small Sync HotStuff RR v4.3}\rhead{\small L1 授权的区域会话}
\cfoot{\thepage}
\setlength{\parindent}{0pt}\setlength{\parskip}{4pt}
\setlength{\emergencystretch}{2em}
\renewcommand{\topfraction}{.92}
\renewcommand{\textfraction}{.06}
\renewcommand{\floatpagefraction}{.85}
\makeatletter
\setlength{\@fptop}{0pt}\setlength{\@fpbot}{0pt plus 1fil}
\makeatother
\setlist{nosep,leftmargin=*}
\newcolumntype{Y}{>{\raggedright\arraybackslash}X}
\definecolor{rrRed}{HTML}{A51C30}
\definecolor{rrBlue}{HTML}{2563EB}
\definecolor{rrTeal}{HTML}{0F766E}
\newcommand{\RR}[1]{{\color{rrRed}#1}}
\newcommand{\QC}{\mathcal C}
\newcommand{\FC}{\mathcal F}
\newcommand{\Hash}{\mathsf H}
\newcommand{\rank}{\operatorname{rank}}
\newcommand{\End}{\operatorname{End}}
\newcommand{\TB}{\mathit{TB}}
\newcommand{\tagg}[1]{\ensuremath{\mathsf{#1}}}
\newtheorem{definition}{定义}[section]
\newaliascnt{assumption}{definition}
\newtheorem{assumption}[assumption]{假设}
\aliascntresetthe{assumption}
\crefname{assumption}{假设}{假设}
\Crefname{assumption}{假设}{假设}
\newcounter{protocol}
% Protocol boxes float to the top of a page, so long boxes leave no gaps.
\newenvironment{protocol}[2]{\begin{figure}[htbp]\refstepcounter{protocol}\label{#1}
\begin{mdframed}[linewidth=.6pt,skipabove=0pt,skipbelow=0pt,
innertopmargin=7pt,innerbottommargin=7pt]\small
\textbf{协议 \theprotocol. #2}\par\smallskip}{\end{mdframed}\end{figure}}
\crefname{protocol}{协议}{协议}
\Crefname{protocol}{协议}{协议}
\title{Sync HotStuff RR v4.3\\\large 区域终局性与 L1 授权轮换}
\author{}\date{2026 年 10 月 12 日}
\crefname{section}{节}{节}
\Crefname{section}{节}{节}
\crefname{table}{表}{表}
\Crefname{table}{表}{表}
\crefname{figure}{图}{图}
\Crefname{figure}{图}{图}
\crefname{definition}{定义}{定义}
\Crefname{definition}{定义}{定义}
\renewcommand{\thesection}{\arabic{section}}
\renewcommand{\thesubsection}{\thesection.\arabic{subsection}}
\crefname{subsection}{节}{节}
\Crefname{subsection}{节}{节}
\begin{document}
\maketitle
\begin{abstract}
本草案将 Sync HotStuff RR v3.3~\cite{rr33} 扩展为多个区域轮流运行的协议。一个区域的一轮服务称为一个\emph{会话}（session）。L1 为每个会话固定质押配置、领导者排程及 DATA 区块配额。在会话内部，该区域运行 v3.3：提案计时器负责提交区块，末位证书负责轮换领导者。会话有两种结束方式。正常交接时，区域提交其配额内的最后一个区块，下一个处于 live 状态的区域从该区块继续。强制交接时，一个诚实的多签委员会在 L1 发布强制指令。该指令指定外部链从此会话保留的最后一个 DATA 高度。旧区域仍可能提交更高的区块，但外部链会丢弃这些区块。两种交接使用相同的引导步骤。若一次恢复从配额上限处开始，则提出一个 DRAIN 区块。它不携带数据，只用于提交最后一个 DATA 区块。本草案规定模型及协议；安全性和复杂度分析留待后续。
\end{abstract}

\section{概览}
\label{sec:overview}
RR v3.3~\cite{rr33} 针对单个区域改写了标准同步模型下的 Sync HotStuff~\cite{sync}。本草案让多个区域轮流生产区块。会话 $r$ 表示一轮服务，由区域 $i_r$ 承担。设计分为三层。
\begin{itemize}
\item \textbf{L1} 是公开的授权来源。它固定每个会话的投票者、领导者和配额，同时承载状态更新及强制指令（\cref{sec:model}）。
\item \textbf{区域引擎} 在一个会话内运行 v3.3，并限制 DATA 区块配额（\cref{sec:regional,sec:protocols}）。其副本是区域 $i_r$ 的成员。
\item \textbf{外部链} 是客户端及后续区域接受的链。每个节点用各会话保留的 DATA 前缀构造该链。下一个区域从它的末端开始（\cref{sec:external,sec:xprotocols}）。这里的节点包括所有区域的副本及任何其他验证者。
\end{itemize}
骨干网络（backbone）在区域间传送区块和回执。它从不选择下一个区域；选择由 L1 记录和区域证书决定。

会话在配额内最后一个区块 $T_r$ 提交后，通过\emph{正常}交接结束；也可以在 L1 强制指令于 DATA 高度 $h_*$ 截断其外部历史时，通过\emph{强制}交接结束。\cref{sec:examples} 中的\cref{fig:endings} 展示了这两种结束方式。

视图为 $(r,g,t)$：会话 $r$、恢复代次 $g$ 和领导者任期 $t$。视图按字典序排序。每个会话从 $(r,0,0)$ 开始。\Cref{tab:views} 列出四种视图变化。前两种与 v3.3 一样，发生在区域 $i_r$ 内部；后两种由下一个区域的成员执行。
\begin{table}[htbp]\centering\small
\caption{四种视图变化。}\label{tab:views}
\begin{tabularx}{\linewidth}{@{}p{.16\linewidth}p{.17\linewidth}Y p{.22\linewidth}@{}}
\toprule
变化 & 新视图 & 触发条件 & 提交计时器\\\midrule
正常轮换 & $(r,g,t+1)$ & 任期 $t$ 的末位 QC，认证的 DATA 高度低于 $H_r$ & 继续运行。\\
恢复 & $(r,g+1,0)$ & 代次 $g$ 的退出证明，随后等待 $\Delta$ & 取消代次 $g$ 的计时器。\\
正常交接 & $(r+1,0,0)$ & 对 REGULAR 见证达到 Ready；见证包含 $T_r$ 的回执（\cref{sec:ready}） & 旧引擎已经完成。\\
强制交接 & $(r+1,0,0)$ & 对 FORCE 见证达到 Ready；见证包含已激活的强制指令 & 旧引擎不变。\\\bottomrule
\end{tabularx}
\end{table}
认证高度 $H_r$ 处 DATA 的 QC，或认证 DRAIN 的 QC，都不启动新任期。不含强制指令的状态更新不改变视图。恢复保留会话原有配额。

\section{模型与假设}\label{sec:model}
\subsection{L1 控制状态}\label{sec:l1}
节点从已经终局化的 L1 读取两类记录（\cref{tab:l1}）。\emph{授权} $\tagg{AUTH}_{r,g}$ 管理代次 $(r,g)$。它包含会话配置、配额及该代次的领导者排程。授权可以由预设的 L1 规则推导，因此一次区域内轮换不需要新的 L1 交易。多签委员会发布\emph{状态更新}。
\begin{table}[htbp]\centering\small
\caption{协议使用的 L1 记录。}\label{tab:l1}
\begin{tabularx}{\linewidth}{@{}p{.27\linewidth}Y@{}}\toprule
字段 & 含义\\\midrule
\multicolumn{2}{@{}l}{\emph{授权} $\tagg{AUTH}_{r,g}$}\\
质押配置 & 各区域的成员 $V_i$、公钥和投票权重，以及 $K$、$\Delta_i$ 等协议参数。\\
领导者排程 & 每个视图的领导者 $L(r,g,t)\in V_{i_r}$。\\
区域配额 $\TB_r$ & 会话 $r$ 可以新增的 DATA 区块数；恢复不重置配额。\\\midrule
\multicolumn{2}{@{}l}{\emph{状态更新}}\\
状态映射 $\chi$ & 每个区域处于 live（可服务）还是 off（停用）状态。\\
强制标志 & 若此更新强制执行交接，则为 true。\\
目标会话 $r$ & 要截断的会话，仅在强制标志为 true 时使用。\\
锚点高度 $h_*$ & 外部链从会话 $r$ 保留的最后一个 DATA 高度，仅在强制标志为 true 时使用。它是高度，而非区块哈希。\\
激活高度 $\lambda_a$ & L1 区块 $\lambda_a$ 终局化后，该更新生效。\\\bottomrule
\end{tabularx}
\end{table}
排程不固定一个会话由哪个区域服务。区域 $i_r$ 在开启会话 $r$ 的交接中选定（\cref{sec:witness}）；随后排程指定 $i_r$ 内部的领导者。

状态更新遵循以下规则。
\begin{itemize}
\item 每次更新都必须在激活高度之前至少提前一个配置好的通知期发布。节点可以在这段时间内准备；更新在激活之前没有效力。
\item 相邻更新的激活高度严格递增。节点按 L1 顺序应用更新。在 L1 高度 $\lambda$ 处的\emph{有效状态映射} $\chi(\lambda)$，是激活高度不超过 $\lambda$ 的最新更新中的映射。
\item 强制更新同时将其目标区域标记为 off。强制标志为 false 的更新只改变状态映射。
\item 节点检查每次更新的授权及 L1 终局性。已经终局化的记录永不改变。
\end{itemize}
初始配置固定区域循环顺序 $\mathcal L_R$、初始状态映射及创世区块。$\mathcal L_R$ 中第一个 live 区域服务会话 $0$。若没有 live 区域，初始化等待。

\subsection{质押与故障}\label{sec:stake}
区域 $i$ 的成员 $v$ 具有投票权重 $\omega_{i,v}>0$。令 $W_i=\sum_{v\in V_i}\omega_{i,v}$。区域 $i$ 中拜占庭成员的总权重至多为 $F_i<W_i/2$。每个区域证书需要来自不同成员、内容一致且总权重严格大于 $F_i$ 的签名。因此，每份证书都包含一个正确签名者，且仅凭正确成员就可以组成证书。这替代了 v3.3 中 $n=2f+1$ 时的 $f+1$ 门槛。

成员和权重在整个会话内保持不变，包括该会话的恢复期间。罚没（slashing）及状态更新不改变它们。后续会话根据旧快照验证旧证书。上界 $F_i$ 计入该快照中所有在其证书仍被接受期间遭到攻破的密钥。在证书不再被接受之后攻破旧密钥，不属于本模型的范围。

正确副本保存其协议状态和已提交区块。签名不可伪造，哈希具有绑定性。

\subsection{假设}\label{sec:assumptions}
\begin{assumption}[L1 授权与控制顺序]\label{ass:authority}\leavevmode\par
\begin{itemize}
\item L1 在区域 $i_r$ 的每个正确成员最早可能进入视图 $(r,g,0)$ 之前，向它提供 $\tagg{AUTH}_{r,g}$。恢复视图也包括在内。
\item 所有正确节点看到相同的、已经终局化的 L1 历史。它们也按一个共同顺序看到 L1 更新和正常交接，该顺序称为\emph{控制顺序}（control order）。
\item L1 独立于区域 BFT 持续推进。
\end{itemize}
\end{assumption}
正常交接不向 L1 发布任何内容。本草案将控制顺序作为一项服务假设，并不规定它如何构建。任何单个节点上的消息到达顺序都不能决定该顺序。

\begin{assumption}[经认证的同步通信与全局时钟]\label{ass:network}\leavevmode\par
\begin{itemize}
\item 消息使用经认证的通信信道。
\item 在区域 $i$ 内，正确成员到正确成员的消息实际延迟为 $\delta\le\Delta_i$；不同消息的 $\delta$ 可以不同。
\item 所有节点共享一个全局时钟。
\end{itemize}
\end{assumption}
延迟上界包含完整验证依赖的传送及处理时间。每个副本在全局时钟上记录自己的进入事件和投票事件。

\begin{assumption}[跨区域骨干网络]\label{ass:backbone}
骨干网络最终会将每个会话的已提交区块及其区域回执交付给所有其他节点。其延迟没有固定上界。
\end{assumption}

\begin{assumption}[诚实的多签委员会]\label{ass:committee}
多签委员会是诚实的。它发布的每次更新都遵循\cref{sec:l1,sec:force} 的规则。
\end{assumption}

\subsection{推进条件}\label{sec:progress}
留待后续的活性分析还将使用以下条件。
\begin{itemize}
\item \textbf{排程。}正确领导者在实际进入的视图中反复获得服务机会。这包括恢复初始化领导者，以及跨会话中生产 DATA 的任期。从未到达的视图中的任期不计入。
\item \textbf{排除。}被揭露存在矛盾签名行为（equivocation）的领导者最终被后续授权排除，与 v3.3 的问责假设一致。
\item \textbf{委员会服务。}委员会给工作正常的区域足够的连续时间来生产 DATA 并完成交接。
\item \textbf{请求。}正确客户端的请求最终到达服务区域及待处理池（\cref{sec:requests}）。它们具有共同的准入顺序，每个请求之前只有有限个请求。每个请求都能装入一个批次。
\end{itemize}
正确领导者向每个 DATA 批次填入最早的、可准入且能够装入批次的待处理请求，并排除已出现在其提案祖先链中的请求。若不存在这样的请求，可以提出空 DATA 区块。

\section{区域引擎：定义}\label{sec:regional}
\subsection{区块、位置与 L1 引用}\label{sec:blocks}
\emph{会话锚点} $A_r$ 是会话 $r$ 从会话 $r-1$ 继承的 DATA 区块；$A_0$ 是创世区块。$A_r$ 已经持有其所属会话的区域回执；创世区块持有一份合成回执。\emph{起始区块}（opening）$A_{r,g}$ 是代次 $g$ 开始的区块，其中 $A_{r,0}=A_r$。对于 $g>0$，NEWVIEW 选择 $A_{r,g}$（\cref{alg:recovery}）。该区块已获认证，但可能尚未提交。它在此代次内永不改变。

每个区块有两个位置量。
\begin{itemize}
\item \emph{DATA 高度} $h(B)$ 统计从创世区块到 $B$ 的链上的 DATA 区块数。
\item \emph{物理位置} $k_r(B)$ 统计会话 $r$ 内的 DATA 和 DRAIN 区块数，且 $k_r(A_r)=0$。这一局部计数不改变 $A_r$ 的已签名字段。
\end{itemize}
会话 $r$ 的新区块形式为
\[
 B=(r,k,h,\mathit{type},x,\Hash(B_{\rm parent}),\rho),
 \qquad \mathit{type}\in\{\tagg{DATA},\tagg{DRAIN}\},
\]
其中 $k=k_r(B_{\rm parent})+1$。DATA 区块携带请求批次 $x$，且 $h=h(B_{\rm parent})+1$。DRAIN 区块不携带批次，且 $h=h(B_{\rm parent})$。

\paragraph{L1 引用。}
引用 $\rho$ 指定一个已经终局化的 L1 区块。它包含该 L1 区块的高度 $\lambda(B)$ 和哈希，以及截至该高度的 L1 控制记录的注册表证明。必须满足以下规则：
\begin{itemize}
\item 沿 DATA 祖先链，$\lambda$ 不递减；
\item 会话 $r$ 的第一个 DATA 区块满足 $\lambda\ge\lambda^{\rm open}_r$，后者为开启会话 $r$ 的边界所对应的选择高度（\cref{sec:witness}）；
\item DRAIN 区块复制 $T_r$ 的引用。
\end{itemize}
$T_r$ 的引用选择正常交接使用的状态映射（\cref{sec:witness}），但不决定该交接在控制顺序中的位置。

区域有效性只使用会话固定的授权。后续 L1 更新不改变该授权。一个区域内有效的区块仍可能被排除在外部链之外（\cref{sec:external}）。

\subsection{配额、任期与 DRAIN}\label{sec:quota}
\paragraph{会话配额。}
会话 $r$ 可以新增 $\TB_r\ge1$ 个 DATA 区块。其末端 DATA 高度为
\[
 H_r=h(A_r)+\TB_r,
\]
$T_r$ 表示高度 $H_r$ 处的 DATA 区块。配额统计提案祖先链中的所有 DATA 区块，无论是否已提交。DRAIN 区块、重发及 ANCHOR 重新认证均不消耗配额。

\paragraph{DATA 代次与收尾代次。}
若 $g>0$ 且 $h(A_{r,g})=H_r$，则代次 $(r,g)$ 是一个\emph{收尾代次}（drain generation）。否则，它是一个\emph{DATA 代次}。代次 $(r,0)$ 始终是 DATA 代次。

\paragraph{DATA 代次的任期。}
当 $h(A_{r,g})+tK<H_r$ 时，任期 $t$ 存在。它有
\[
 m_{r,g,t}=\min\{K,\;H_r-h(A_{r,g})-tK\}
\]
个新位置（fresh positions），$j=0,\ldots,m_{r,g,t}-1$。位置 $j$ 放置一个 DATA 区块，其位置量为
\[
 h=h(A_{r,g})+tK+j+1,\qquad k=k_r(A_{r,g})+tK+j+1.
\]
任期的\emph{末位}是 $j=m_{r,g,t}-1$。只有最后一个任期可以少于 $K$ 个位置。高度 $H_r$ 处的末位区块之后不存在新任期。

例如，令 $h(A_r)=100$、$\TB_r=10$、$K=4$。三个任期依次生产 DATA 101--104、105--108 和 109--110。若恢复改为从 DATA 107 开始，则恢复后的第一个任期有三个位置：108--110。恢复永不改变 $H_r=110$。

\paragraph{收尾代次。}
一个收尾代次只有一个新位置 $(t,j)=(0,0)$。它也是末位，且 $m_{r,g,0}=1$。该位置的区块是 DRAIN，其父区块是 $A_{r,g}$，以本代次的 ANCHOR QC 为依据。不存在其他位置，也不存在下一任期。

DRAIN 的存在只有一个原因。与 v3.3 一样，ANCHOR QC 不启动计时器。因此，恢复后的起始区块只有在一个新子区块的计时器到期时才会提交。若起始区块处于 $H_r$，配额禁止再添加 DATA 子区块。DRAIN 子区块提供所需计时器。提交它会同时提交其祖先，包括 $T_r$。另外还有三条规则。
\begin{itemize}
\item 只有收尾代次提出 DRAIN。到达 $H_r$ 的 DATA 代次不能再追加 DRAIN；它等待 $T_r$ 的计时器，或在退出后恢复。
\item 收尾代次的起始区块可以是先前的 DRAIN。每次被中断的恢复都可能增加一个 DRAIN，因此 DRAIN 的数量没有上界。
\item 此规则只取决于 $h(A_{r,g})$，不取决于某个副本是否已经提交 $T_r$。
\end{itemize}
DRAIN 区块保留在区域日志中。外部链从不包含它们，下一会话从 $T_r$ 开始。

\emph{新位置}是上面定义的一个位置 $(t,j)$。\emph{新区块}是某个新位置处的 DATA 或 DRAIN 区块。起始区块不属于新区块。\emph{末位 QC} 认证一个末位上的区块。

\subsection{证书与签名消息}\label{sec:certificates}
\Cref{tab:certificates} 列出一个会话的证书。每种证书都需要来自 $V_{i_r}$ 中不同成员、总权重 $>F_{i_r}$ 的一致签名。
\begin{table}[htbp]\centering\small
\caption{一个会话的证书。}\label{tab:certificates}
\begin{tabularx}{\linewidth}{@{}p{.25\linewidth}p{.24\linewidth}Y@{}}\toprule
证书 & 签名消息 & 用途\\\midrule
提案 QC $\QC_{r,g,t}(B)$ & DATA-VOTE 或 DRAIN-VOTE & 为子区块提供依据。低于 $H_r$ 的末位 QC 还准入下一任期。\\
ANCHOR QC & ANCHOR-VOTE & 在提出任何新区块之前，于代次 $g$ 内重新认证起始区块 $A_{r,g}$。\\
BLAME 证明 & BLAME & 冻结一个代次（\cref{sec:quit}）。\\
区域回执 $\FC_r(B)$ & FINAL & 证明一个正确副本已提交 DATA 区块 $B$。\\\bottomrule
\end{tabularx}
\end{table}
正确副本对每个 $(r,g,t,j)$ 至多投一次新区块票，对每个 $(r,g)$ 至多投一次 ANCHOR 票。DATA 和 DRAIN 投票遵循相同规则，但属于不同签名类型。任何 QC 和任何 ANCHOR 投票都不启动提交计时器。

\paragraph{签名上下文。}
每个签名都绑定链、投票配置、会话 $r$、区域 $i_r$ 及消息类型。视图内的消息还绑定 L1 授权及 $(r,g,t)$。其他签名字段如下：
\begin{itemize}
\item \textbf{PROPOSE 与新区块投票：}起始区块、位置 $j$、区块类型、物理位置、DATA 高度及区块哈希；
\item \textbf{NEWVIEW 与 ANCHOR-VOTE：}所选起始区块的身份。NEWVIEW 还绑定其依据和进入证明；
\item \textbf{STATUS：}发送者的进入锁及进入证明；无需指定起始区块；
\item \textbf{BLAME：}只绑定视图，不绑定起始区块、区块哈希或进度计数。因此，同一视图的所有抱怨都签署相同声明，哪怕 NEWVIEW 尚未出现；
\item \textbf{FINAL：}仅绑定在会话内稳定的字段，见下文。
\end{itemize}
消息展示式省略这些上下文字段。

\paragraph{等级与基础依据。}
在会话 $r$ 内，于任期 $(g,t)$ 为区块 $B$ 形成的 QC，等级为 $(g,t,k_r(B))$，按字典序比较。代次 $g$ 的 ANCHOR QC 的等级为 $(g,0,k_r(A_{r,g}))$，恰好低于该代次的第一个新区块 QC。对已认证内容采用固定全序以打破平局；签名者集合及编码不影响等级。起始区块的身份为 $(g,k_r(A_{r,g}),h(A_{r,g}),\Hash(A_{r,g}))$。

锚点 $A_r$ 还具有一个\emph{基础依据} $Q_r^0$：即开启会话 $r$ 的边界见证中 $A_r$ 的回执（\cref{sec:witness}）。$Q_r^0$ 的等级低于会话 $r$ 的所有 QC。当恢复时不知道本会话的任何 QC，就使用它。另一会话的 QC 从不作为会话 $r$ 中的依据，因此不能替换 $A_r$。

\paragraph{区域回执。}
副本在通过 DATA 区块 $B$ 自己的计时器提交它，或将其作为祖先提交后，向所在区域及骨干网络发送
\[
 \langle\tagg{FINAL},\mathit{cfg}_r,r,h(B),\Hash(B)\rangle_v.
\]
总权重 $>F_{i_r}$ 的一致 FINAL 消息形成区域回执 $\FC_r(B)$。回执携带 $B$ 及其所需祖先，或引用一个已经验证的缓存。
\begin{itemize}
\item FINAL 不绑定代次或任期，因此，同一区块的 FINAL 消息可以跨恢复合并。区块哈希保留了区块的签名上下文。
\item 若 DRAIN 提交 $T_r$，副本为 $T_r$ 发送 FINAL，从不为 DRAIN 发送。
\item 会话 $r$ 的副本不为 $A_r$ 发送 FINAL；$A_r$ 保留会话 $r-1$ 的回执。
\item 副本可以提交其自身会话的一份有效回执所对应的 DATA 前缀。其他会话的回执只充当历史或锚点证据，从不改变本会话的锁或日志。
\end{itemize}

\subsection{准入、验证与引擎状态}\label{sec:admission}
\paragraph{准入。}
副本 $v$ 按以下规则准入视图。
\begin{itemize}
\item $(r,0,0)$ 在 $v$ 对开启会话 $r$ 的边界达到 Ready 时准入（\cref{def:ready}）。见证是进入证明。NEWVIEW 必须恰好选择 $A_r$。
\item 对 $g>0$，当 $v$ 持有 $\tagg{AUTH}_{r,g}$ 及代次 $g-1$ 的退出证明时，$(r,g,0)$ 准入。该退出证明就是进入证明。NEWVIEW 以会话 $r$ 的较早代次的 QC 或 $Q_r^0$ 为依据。
\item 对 $t>0$，当 $v$ 还持有任期 $0,\ldots,t-1$ 的末位 QC，且每份 QC 均认证低于 $H_r$ 的 DATA 时，$(r,g,t)$ 准入。
\end{itemize}
第零任期在其 NEWVIEW 到达前就已准入。因此，即使初始领导者保持沉默，进度检查也会运行。新区块提案需要本代次的 ANCHOR QC。此规则适用于所有代次，包括会话 $0$ 的代次 $0$；与 v3.3 不同，创世区块没有捷径。若消息所属视图已经准入，且签名有效，则消息是\emph{可准入的}。

\paragraph{完整验证。}
验证检查签名、授权、准入、区块类型、配额、两个位置量、L1 引用规则、父区块规则（\cref{tab:parent}）、父 QC、哈希链接，以及所需祖先链中各区块的应用有效性。每个相关祖先都必须包含其签名提案及区块数据；只有一串证书哈希并不够。依赖项来自消息或已经验证的本地缓存。正确发送者发送完整验证包。只有当接收者能在\cref{ass:network} 的延迟上界内解析引用时，才允许使用该引用。验证包中每项依赖只列一次。

\paragraph{证据优先。}
在验证消息或丢弃重复消息之前，检查其中所有已签名材料是否构成退出证明（\cref{sec:quit}）。证明本身需要有效签名及准入见证。即使携带它的消息无效或不完整，有效证明仍然有效。应将收到的材料与所有已保留材料比较，而不只比较该消息内部的对象。

\paragraph{统一的提案处理规则。}
当前代次中每个已经准入、且 $v$ 尚未为其投票的新区块提案，都调用 Vote 步骤。这包括附带历史中的提案、较早任期的提案，以及区块已作为祖先提交的提案。提案身份为 $(r,g,t,j,\Hash(B))$。父 QC 的签名者集合或编码不改变此身份。在检查证据之后，按此身份去重。旧代次提案仅保留，不产生新投票或计时器。未来代次消息等待进入该代次后处理。

\paragraph{引擎状态。}
副本对其所在区域的每个会话运行一个引擎。会话 $r$ 的引擎处于三种状态之一。
\begin{itemize}
\item \textbf{Active（活动）：}当前代次已经准入且未被冻结。引擎执行提案、投票和任期变化。
\item \textbf{Frozen（冻结）：}当前代次存在退出证明。引擎停止提案、投票和任期变化，但仍处理证据、旧投票、QC、回执及恢复。
\item \textbf{Complete（完成）：}副本持有有效回执 $\FC_r(T_r)$。会话 $r$ 的生产结束，即使外部交接仍在等待。引擎只继续提供数据。
\end{itemize}
本地提交 $T_r$ 不会使引擎完成。在回执形成之前，其他成员仍可能需要该副本参与恢复投票。每条消息都交给其所属会话的引擎，因此，一个会话的证据从不冻结另一个会话。这些状态独立于\cref{sec:states} 的外部状态。外部链可能已进入后续会话，而旧引擎仍在运行。
\begin{table}[htbp]\centering\small
\caption{任期 $(r,g,t)$ 的下一个新位置 $j$ 所使用的父区块选择。}\label{tab:parent}
\begin{tabularx}{\linewidth}{@{}l Y Y@{}}\toprule
条件 & $B_{\mathrm{parent}}$ & $Q_{\mathrm{parent}}$\\\midrule
$j=0,t=0$（DATA 或 DRAIN） & 起始区块 $A_{r,g}$ & 它在 $(r,g,0)$ 中的 ANCHOR QC\\
$j=0,t>0$ & $(r,g,t-1)$ 的末位区块 & 它在 $(r,g,t-1)$ 中的提案 QC\\
$1\le j<m_{r,g,t}$ & $(r,g,t)$ 的位置 $j-1$ & 它在 $(r,g,t)$ 中的提案 QC\\\bottomrule
\end{tabularx}
\end{table}

\section{区域协议}\label{sec:protocols}
协议 1 和 2 在一个会话 $r$ 的引擎内运行。它们是在 v3.3 的协议上应用以下改动得到的。\RR{红色文字}标记协议框、消息处理顺序及退出证明中相对 v3.3 的每项差异。说明文字放在协议框外。

\paragraph{记号。}
协议框内的任期 $(g,t)$ 表示视图 $(r,g,t)$。消息省略 $r$，但签名上下文绑定了它。$L(g,t)$、$e_v(g,t)$ 和 $m_{g,t}$ 分别代表 $L(r,g,t)$、$e_v(r,g,t)$ 和 $m_{r,g,t}$。$\Delta$ 和 $F$ 分别代表 $\Delta_{i_r}$ 和 $F_{i_r}$。v3.3 中的副本 $r$ 在这里记为副本 $v$，其区块 $B_k$ 记为 $B$，其检查索引 $i$ 记为 $p$。这些更名不标红。

\paragraph{相对 v3.3 的改动。}
\begin{enumerate}[label=\arabic*.]
\item \emph{质押。}证书需要权重 $>F$，而非 $f+1$ 个签名。
\item \emph{配额。}一个任期有 $m_{g,t}\le K$ 个位置。$H_r$ 之后不存在新任期。
\item \emph{DRAIN。}收尾代次提出一个 DRAIN。对它的投票计入进度。
\item \emph{回执。}副本提交 DATA 后发送 FINAL。$T_r$ 的回执使引擎完成（\cref{alg:steady} 新增的第 5 步）。
\item \emph{计时器。}完成也会取消计时器。强制指令从不取消计时器。
\item \emph{恢复进入。}副本在 $a_v+\Delta$ 进入，因为授权会及时到达（\cref{ass:authority}）。$Q_r^0$ 可以充当进入锁及依据。
\item \emph{会话。}每条消息交给其所属会话的引擎。会话 $0$ 通过引导开启（\cref{alg:bootstrap}）。
\end{enumerate}

\paragraph{消息处理顺序。}\leavevmode
\RR{副本 $v$ 首先将每条消息交给其签名所绑定的会话引擎。L1 更新、见证及其他会话的回执交给\cref{alg:external}。}引擎按以下顺序处理每条收到的消息：
\begin{enumerate}[label=(\arabic*)]
\item 检查所有已签名材料是否构成退出证明；
\item 检查完整性及有效性；
\item 按父区块优先的顺序，在 Vote 步骤下处理附带的 DATA \RR{和 DRAIN} 提案；
\item 按任期递增顺序，在\cref{alg:steady} 第 4 步下处理附带的末位 QC；
\item 处理消息本身。
\end{enumerate}
当一条消息与一个计时器在同一时刻触发时，$v$ 先处理消息。

\subsection{稳态}\label{sec:steady}
\Cref{alg:steady} 给出会话的稳态流程：提案、投票、计时器、任期轮换和完成。
\begin{protocol}{alg:steady}{提案、投票、提交、正常轮换\RR{及完成}}
副本 $v$ \RR{在会话 $r$ 的引擎内}运行以下步骤。当前任期为 $(g,t)$。
\begin{enumerate}[label=\textbf{\arabic*)},itemsep=3pt]
\item \textbf{Propose（提案）。}仅当 $v=L(g,t)$、$g$ 未冻结\RR{且会话 $r$ 尚未完成}时运行。
\RR{在 DATA 代次中，}按顺序对每个位置 $j=0,\ldots,\RR{m_{g,t}}-1$：
\begin{itemize}
\item 等待 $v$ 持有位置 $j$ 所需的 $B_{\mathrm{parent}}$ 和 $Q_{\mathrm{parent}}$（\cref{tab:parent}）。
\item 构造\RR{$B=(r,k,h,\tagg{DATA},x,\Hash(B_{\mathrm{parent}}),\rho)$，其中 $k$、$h$ 遵循\cref{sec:quota}，$\rho$ 为 L1 引用}。批次 $x$ 遵循\cref{sec:progress}。
\item 广播 $P=\langle\tagg{PROPOSE},g,t,j,B,Q_{\mathrm{parent}}\rangle_v$ 及其验证包。
\item 按 Vote 处理 $P$。
\end{itemize}
\RR{在收尾代次中，只在 $(t,j)=(0,0)$ 提案一次。区块为 DRAIN，父区块为 $A_{r,g}$，不携带批次，DATA 高度为 $H_r$，L1 引用为 $T_r$ 的引用。按上述方式广播并处理它。}
不等待父区块提交。从不为同一位置签署两个不同区块。重发时重复相同的 $P$。

\item \textbf{Vote（投票）。}对任期 $(g_P,t_P)$ 的位置 $j_P$ 上每个 DATA \RR{或 DRAIN} 提案 $P$ 运行本步骤。这里 $t_P$ 可以早于 $t$。本步骤也适用于附带历史内的提案。
\begin{enumerate}[label=(\alph*)]
\item 若 $P$ 或其附带材料包含 $g$ 的退出证明，运行 Quit 并停止。
\item 若 $g_P<g$，保留 $P$ 并停止。若 $g_P>g$，缓存 $P$，直到 $v$ 进入代次 $g_P$。
\item 若 $g$ 已冻结\RR{或会话 $r$ 已完成}，保留 $P$ 并停止。
\item 若缺少依赖项，保留 $P$、请求依赖项并停止。依赖项到达后重跑本步骤。
\item 若 $P$ 无效，丢弃并停止。\RR{有效性包括区块类型、配额、两个位置量及 L1 引用。}
\item 若 $v$ 已经为 $P$ 投票，停止。重复消息从不重新启动计时器。
\item 否则，令 $\tau_v(P)$ 为当前时间。向所有副本转发 $P$ 及其验证包。广播
$\langle\tagg{DATA\mbox{-}VOTE},g,t_P,j_P,\Hash(B)\rangle_v$
\RR{（若 $B$ 为 DRAIN，则使用 DRAIN-VOTE）}。启动 $P$ 的提交计时器，到期时刻为 $\tau_v(P)+2\Delta$。
\end{enumerate}

\item \textbf{Commit（提交）。}当 $P$ 的计时器到期，提交 $B$ 及其尚未提交的祖先。\RR{对这次提交的、属于会话 $r$ 的每个 DATA 区块 $D$，向本区域及骨干网络发送
$\langle\tagg{FINAL},\mathit{cfg}_r,r,h(D),\Hash(D)\rangle_v$。
提交 DRAIN 时，为 $T_r$ 产生 FINAL，从不为 DRAIN 本身产生 FINAL。将 $v$ 获得的每份会话 $r$ 的回执转发给本区域及骨干网络。}
只有 Quit \RR{和 Completion} 取消计时器。进入新任期不会取消它。QC 从不启动计时器。\RR{强制指令从不取消计时器。}

\item \textbf{末位 QC 与轮换。}当 $v$ 首次获得任期 $(g,t')$ 的有效末位 QC 时运行本步骤。“获得”表示 $v$ 从投票组装出该 QC，或在任意消息中发现它，包括作为另一个提案的 $Q_{\mathrm{parent}}$ 发现它。
\begin{enumerate}[label=(\alph*)]
\item 按父区块优先顺序，在 Vote 下处理该 QC 历史中每个尚未处理的提案。
\item 若 $g$ 此时已经冻结，停止。\RR{若 QC 认证的是 $H_r$ 处的 DATA 或一个 DRAIN，停止：不存在下一任期。}
\item 向 $L(g,t'+1)$ 发送该 QC 及其验证包。
\item 若 $t\le t'$，进入 $(g,t'+1)$，并将 $e_v(g,t'+1)$ 设为当前时间。取消所有较早任期的进度检查。
\end{enumerate}
QC 从不让 $v$ 回到较早任期，也从不改变已记录的进入时间。

\item \RR{\textbf{Completion（完成）。}当 $v$ 首次获得有效区域回执 $\FC_r(T_r)$ 时运行本步骤。}
\begin{itemize}\color{rrRed}
\item 提交 $T_r$ 及其尚未提交的祖先。
\item 将会话 $r$ 标记为完成。停止其提案、投票、进度检查、任期变化及恢复。取消其尚未到期的计时器。
\item 向本区域及骨干网络转发 $\FC_r(T_r)$ 及其验证包。
\item 继续提供数据。交接由\cref{alg:external} 执行。
\end{itemize}
\RR{强制指令从不运行本步骤。}
\end{enumerate}
\end{protocol}

接受时间 $\tau_v(P)$ 是完成完整验证并首次投票的时刻，不是不完整数据包较早到达的时刻。每个副本记录自己进入每个任期的时间。在会话 $r$ 中，$e_v(0,0)$ 在\cref{alg:bootstrap} 第 2 步、达到 Ready 时设置。对于 $g>0$，$e_v(g,0)$ 在\cref{alg:recovery} 第 3 步设置。对于 $t>0$，$e_v(g,t)$ 在\cref{alg:steady} 第 4 步设置。因此，一个任期的进度检查在该任期本地准入时启动，而不是在其领导者的首个提案到达时启动。沉默的新领导者不能推迟这些检查。

高度 $H_r$ 处的提案或 QC 不会结束会话。若没有退出，$T_r$ 的计时器会提交它，无需更后面的区块。只有 $T_r$ 的回执才会使引擎完成。

\subsection{退出证明与恢复}\label{sec:quit}
退出证明冻结一个代次。随后\cref{alg:recovery} 将会话推进到下一个代次。

\paragraph{完整的退出证明。}
对于\RR{会话 $r$ 中}已经准入的代次 $g$，证明由以下任一种材料组成：
\begin{enumerate}[label=(\alph*)]
\item 对已准入任期的同一条 $\langle\tagg{BLAME},g,t\rangle$ 消息的不同签名\RR{，总权重 $>F$}，连同该任期的准入见证；
\item 由 $L(g,t)$ 在同一个 $(g,t,j)$ 上签署的两个不同的、可准入的 DATA \RR{或 DRAIN} 区块，连同它们的准入\RR{及配额}见证；或者
\item 由 $L(g,0)$ 签署的两份有效 NEWVIEW，为代次 $g$ 选择了不同起始区块。
\end{enumerate}
\RR{超出配额的 DATA 区块，或在收尾代次之外提出的 DRAIN，均无效。它们会被拒绝，且从不计入 (b)。}
即使两个起始区块位于同一条链上、可以按祖先关系比较，选择不同起始区块仍构成矛盾签名，因为它们固定了不同的起始高度 $h(A_{r,g})$。此时高度公式将同一位置映射到不同高度。

同一起始区块使用不同 QC 不构成矛盾签名。来自任意已准入任期的证据都可以冻结其代次，包括本地尚未到达的任期。旧代次证明归档保存，并可提交给 L1。\RR{来自其他会话的证据从不冻结本引擎。}未来代次证明缓存至进入该代次后处理。两者均不重置无关的恢复。

\begin{protocol}{alg:recovery}{抱怨、退出与恢复}
副本 $v$ \RR{在会话 $r$ 的引擎内}运行以下步骤。当前任期为 $(g,t)$。
\begin{enumerate}[label=\textbf{\arabic*)},itemsep=3pt]
\item \textbf{Blame（抱怨）。}\RR{在 DATA 代次中令 $q=m_{g,t}$，在收尾代次中令 $q=1$。}对每个 $p=1,\ldots,\RR{q}$，在时刻
\[
 e_v(g,t)+(2p+4)\Delta
\]
运行一次进度检查。
\RR{进入时即安排 $p=1$ 的检查，即使尚不知道起始区块。知道 $q$ 后再安排其他检查。若某次检查的预定时间已经过去，则立即运行。}
到该时刻，若以下条件全部成立，则广播 $\langle\tagg{BLAME},g,t\rangle_v$：
\begin{itemize}
\item $v$ 当前任期仍为 $(g,t)$，$g$ 未冻结\RR{且会话 $r$ 尚未完成}；
\item $v$ 为任期 $(g,t)$ 的位置投出的 DATA \RR{或 DRAIN} 票数少于 $p$（ANCHOR 票不计入）；
\item $v$ 尚未为 $(g,t)$ 发送 BLAME。
\end{itemize}

\item \textbf{Quit（退出）。}当 $v$ 首次持有其当前代次 $g$ 的退出证明时，运行本步骤。
\begin{itemize}
\item 冻结 $g$。停止 $g$ 中所有进度检查、提案、投票和任期变化。取消 $g$ 的每个尚未到期的提交计时器。
\item 将 $a_v$ 设为当前时间。广播证明。
\item 保留所有提交结果。继续收集 $g$ 的投票及 QC。它们可以提高第 3 步选择的进入锁。\RR{继续处理区域回执（\cref{alg:steady} 第 5 步）。}
\item 若证明属于类型 (b) 或 (c)，则提交给 L1，不等待 L1 处理它。
\end{itemize}

\item \textbf{Status 与进入。}\RR{若会话 $r$ 已完成，跳过本步骤。}在时刻\RR{$e_v(g+1,0)=a_v+\Delta$}，从 L1 读取 $\tagg{AUTH}_{r,g+1}$\RR{（\cref{ass:authority} 保证它已经可用）}。然后：
\begin{itemize}
\item 将进入锁 $Q_v$ 设为 $v$ 已知的、来自代次 $0,\ldots,g$ 的最高等级 QC（DATA\RR{、DRAIN} 或 ANCHOR）\RR{；若一个也不知道，则设为 $Q_r^0$}。此后 $Q_v$ 永不改变。稍后到达的 QC 不更新它。
\item 以未冻结状态进入 $(g+1,0)$，并启动其进度检查。\RR{保留会话 $r$ 的 $A_r$、$H_r$ 和投票快照。}
\item 向 $L(g+1,0)$ 发送 $\langle\tagg{STATUS},g+1,0,Q_v,E\rangle_v$ 及其验证包。这里 $E$ 是 $g$ 的退出证明，也是 $g+1$ 的进入证明。
\item 处理为 $g+1$ 缓存的证明。
\end{itemize}

\item \textbf{New-view（新视图）。}仅在 $L(g+1,0)$ 运行本步骤。自它自己的进入时间起等待 $2\Delta$。若 $g+1$ 仍是当前代次且未冻结：
\begin{itemize}
\item 令 $Q^*$ 为领导者已知的、来自代次 $0,\ldots,g$ 的最高等级 QC，包括 STATUS 消息中的 QC\RR{；若一个也不知道，则令 $Q^*=Q_r^0$}。
\item 令 $A_{r,g+1}$ 为 $Q^*$ 认证的区块。\RR{若 $h(A_{r,g+1})<H_r$，则 $g+1$ 为 DATA 代次；否则为收尾代次。}
\item 广播 $\langle\tagg{NEWVIEW},g+1,0,A_{r,g+1},Q^*,E\rangle_{L(g+1,0)}$ 及其验证包。为 $g+1$ 只选择一个起始区块。
\end{itemize}

\item \textbf{Anchor vote（起始区块投票）。}收到 $g+1$ 的 NEWVIEW 时：
\begin{enumerate}[label=(\alph*)]
\item 若其中包含 $g+1$ 的退出证明，运行 Quit 并停止。
\item 验证其领导者、进入证明 $E$、起始区块 $A_{r,g+1}$、$Q^*$ 及依赖项。若缺少数据，请求数据并等待，不投票。若消息无效，丢弃它。
\item 只有以下条件全部成立时才投票：$g+1$ 是当前代次且未冻结；$v$ 尚未在 $g+1$ 投过 ANCHOR 票；且 $\rank(Q^*)\ge\rank(Q_v)$。
\item 投票时，向所有副本转发 NEWVIEW，并广播 $\langle\tagg{ANCHOR\mbox{-}VOTE},g+1,0,\Hash(A_{r,g+1})\rangle_v$。不启动计时器，也不使用新位置。
\end{enumerate}
\RR{总权重 $>F$ 的}一致 ANCHOR 票形成 ANCHOR QC。$L(g+1,0)$ 获得它后，从 $j=0$ 开始运行\cref{alg:steady} 的 Propose\RR{；在收尾代次中，这个区块为 DRAIN}。所有副本，包括未投 ANCHOR 票的副本，都按 Vote 处理由有效 ANCHOR QC 支持的 DATA \RR{和 DRAIN} 提案。
\end{enumerate}
\end{protocol}

进入锁测试只约束 ANCHOR 投票。本地已经提交 $T_r$ 的副本仍继续参与恢复，直到持有 $\FC_r(T_r)$，因为其他成员可能需要它的投票。

\section{外部链与交接}\label{sec:external}
\subsection{会话状态与外部链}\label{sec:states}
每个节点将每个会话维持在以下三种外部状态之一（\cref{fig:states}）。
\begin{itemize}
\item \textbf{Open（开放）：}会话终点尚未固定。其持有有效回执、且高度不超过 $H_r$ 的 DATA 区块是\emph{暂定}输出，即使区域已经提交了它们。
\item \textbf{Truncated（已截断）：}已激活的强制指令固定了终点 $\End(r)=h_*$（\cref{sec:force}）。锚点数据或接任区域可能仍然缺失。$\End(r)$ 此后永不改变。
\item \textbf{Sealed（已封定）：}锚点 $A_{r+1}$、接任区域 $i_{r+1}$ 和选择高度 $\lambda_s$ 已经固定（\cref{sec:witness}）。保留前缀成为最终结果。
\end{itemize}
正常交接封定会话时，$\End(r)=H_r$ 且 $A_{r+1}=T_r$。强制交接封定会话时，$\End(r)=h_*$。其锚点 $A_{r+1}$ 为会话 $r$ 在高度 $h_*$ 处的 DATA 区块；若 $h_*=h(A_r)$，则为 $A_r$ 本身。会话不会返回较早状态，后续状态更新也不会重新开放它。\emph{当前外部会话}是第一个尚未封定的会话。
\begin{figure}[htbp]\centering
% External states of session r (Section 5.1).
% Open -> Sealed (regular handoff); Open -> Truncated -> Sealed (forced handoff).
% No state is ever left for an earlier one.
\begin{tikzpicture}[x=1cm,y=1cm,font=\fontsize{9}{10.5}\selectfont,
  st/.style={draw,rounded corners=3pt,minimum width=2.9cm,minimum height=1cm,
    align=center,inner sep=3pt},
  tr/.style={-{Stealth[length=1.8mm]},semithick},
  lbl/.style={font=\fontsize{8.5}{10}\selectfont,align=center,inner sep=2pt}]
  \node[st,draw=black!60,fill=black!4] (open) at (0,0)
    {\textbf{Open（开放）}\\[-1pt]{\fontsize{8}{9}\selectfont 输出为暂定结果}};
  \node[st,draw=rrRed,fill=rrRed!6] (trunc) at (4.7,-1.9)
    {\textbf{Truncated（已截断）}\\[-1pt]{\fontsize{8}{9}\selectfont $\End(r)=h_*$ 已固定}};
  \node[st,draw=rrTeal,fill=rrTeal!10] (sealed) at (11.2,0)
    {\textbf{Sealed（已封定）}\\[-1pt]{\fontsize{8}{9}\selectfont 锚点与接任区域已固定}};
  \draw[tr] (open) -- node[lbl,above] {正常交接：$\FC_r(T_r)$ 与 live 接任区域；$\End(r)=H_r$} (sealed);
  \draw[tr] (open.south) |- node[lbl,pos=.74,above] {强制指令激活} (trunc.west);
  \draw[tr] (trunc.east) -| node[lbl,pos=.26,above] {锚点数据与 live 接任区域} (sealed.south);
\end{tikzpicture}

\caption{会话 $r$ 的外部状态。正常交接从 Open 进入 Sealed；强制交接经过 Truncated。会话 $r$ 的区域引擎拥有独立状态（\cref{sec:admission}）。}\label{fig:states}
\end{figure}

\paragraph{外部接受规则。}
设会话 $r$ 已截断或已封定。会话 $r$ 的 DATA 区块 $B$ 属于外部链，当且仅当：
\begin{enumerate}[label=(\roman*)]
\item 它持有有效区域回执 $\FC_r(B)$；
\item $h(A_r)<h(B)\le\End(r)$；且
\item $B$ 是 $\End(r)$ 处的区块，或是它的某个 DATA 祖先，且所在链延伸自 $A_r$。
\end{enumerate}
外部链由创世区块及按会话顺序排列的这些区块构成。区块自身的 L1 引用从不覆盖这些规则。

\subsection{强制指令}\label{sec:force}
强制更新的强制标志为 true，包含目标会话 $r$、锚点高度 $h_*$ 和激活高度 $\lambda_a$。诚实委员会仅在以下条件全部成立时发布它：
\begin{itemize}
\item 发布更新时，会话 $r$ 是当前外部会话；
\item $h(A_r)\le h_*\le H_r$，从而保留继承前缀；
\item 会话 $r$ 在高度 $h_*$ 处的 DATA 区块已在区域内提交，且位于该区域工作正常时生产的前缀中。其回执和祖先链已在区域外可用。委员会选择满足这些条件的最高区块。
\end{itemize}
更新指定高度，而不是区块哈希。回执识别具体区块，因为区域安全性只允许每个高度有一个已提交 DATA 区块。

通知期允许节点提前准备。激活安装截断点：一旦 L1 区块 $\lambda_a$ 终局化，$\End(r)=h_*$。即使锚点数据或接任区域仍缺失，也成立。

强制指令只作用于外部链。它不改变旧区域的计时器、锁、配额或已提交日志，也不是退出证明。旧区域仍可能提出并提交高于 $h_*$ 的区块。其回执仍是有效的区域证据，但这些区块永不返回外部链。强制指令不需要旧区域已经停止的证明，也不需要更高区块不存在的证明。

\subsection{边界见证与接任区域}\label{sec:witness}
\paragraph{接任区域。}
令 $\tagg{NextLive}(i,\chi)$ 为 $\mathcal L_R$ 中位于 $i$ 之后、不同于 $i$ 且在 $\chi$ 中为 live 的第一个区域。搜索至多覆盖一个循环。若不存在这样的区域，则 $\tagg{NextLive}(i,\chi)=\bot$。令
\[
 \lambda_0=\begin{cases}
 \lambda(T_r) & \text{正常交接时，}\\
 \max\{\lambda_a,\lambda(A_{r+1})\} & \text{强制交接时。}
 \end{cases}
\]
\emph{选择高度} $\lambda_s$ 为满足 $\lambda\ge\lambda_0$ 且 $\tagg{NextLive}(i_r,\chi(\lambda))\ne\bot$ 的最小 L1 高度。接任区域为 $i_{r+1}=\tagg{NextLive}(i_r,\chi(\lambda_s))$，并令 $\lambda^{\rm open}_{r+1}=\lambda_s$。因此，每个节点计算出相同的接任区域。若 $\lambda_0$ 处没有其他 live 区域，边界等待第一个使其他区域变为 live 的后续状态映射。一旦选定，接任区域永不改变。对于会话 $0$，$\lambda^{\rm open}_0$ 为初始 L1 高度。

\paragraph{哪个事件决定边界。}
节点按控制顺序处理事件（\cref{ass:authority}）。会话 $r$ 的边界由以下事件中先发生者决定：
\begin{itemize}
\item 目标为 $r$ 的强制指令生效；或者
\item 回执 $\FC_r(T_r)$ 进入控制顺序，且存在 live 接任区域，因此正常交接封定会话。
\end{itemize}
在控制顺序的同一位置，强制指令优先。强制指令之后到达的回执不能撤销它。若强制指令的目标已封定，则只改变状态映射。$T_r$ 的 L1 引用选择状态映射，不会把正常交接提前到控制顺序中的更早位置。

\paragraph{见证。}
边界见证的形式为
\[
 \mathcal W=(\mathit{mode},r,i_{r+1},A_{r+1},\FC(A_{r+1}),\lambda_s,\Xi),
 \qquad \mathit{mode}\in\{\tagg{REGULAR},\tagg{FORCE}\}.
\]
证明 $\Xi$ 认证截至该边界的控制顺序、$\lambda_s$ 处的状态映射、任何已激活的强制指令，以及 $\tagg{AUTH}_{r+1,0}$。
\begin{itemize}
\item \textbf{REGULAR：}$A_{r+1}=T_r$，且其祖先链中 $A_r$ 之后恰有 $\TB_r$ 个 DATA 区块。
\item \textbf{FORCE：}$A_{r+1}$ 为高度 $h_*$ 处的区块，并附有回执。若 $h_*=h(A_r)$，使用 $A_r$ 的旧回执，故障区域不签署任何新内容。
\end{itemize}
证书编码不改变边界身份。任何节点都可以检查见证并据此封定会话。只有 $i_{r+1}$ 的成员用它启动引擎。会话 $0$ 使用合成见证：$A_0$ 为带有合成回执的创世区块，$i_0$ 为 $\mathcal L_R$ 中第一个 live 区域。

\subsection{Ready 事件}\label{sec:ready}
\begin{definition}[Ready]\label{def:ready}
设 $v$ 为区域 $i'$ 的成员。$\tagg{Ready}_v(\mathcal W)$ 在 $v$ 首次完整验证以下全部内容时发生：
\begin{enumerate}[label=(\roman*)]
\item 一份已封定见证 $\mathcal W$，指定 $i'$ 为会话 $r$ 的接任区域；
\item 锚点及 $v$ 缺失的任何应用历史；以及
\item 已经终局化的 $\tagg{AUTH}_{r+1,0}$。
\end{enumerate}
REGULAR 见证必须展示 $\FC_r(T_r)$ 及完整配额。FORCE 见证则必须展示截至 $\lambda_a$ 的已终局化 L1，以及控制顺序中的截断点。若 $v$ 已经进入会话 $r+1$ 或更晚会话，则不再发生 Ready。不完整见证持续等待，直到其数据和授权到达。
\end{definition}
Ready 是唯一允许一个区域启动下一会话的事件。高度 $H_r$ 处的区块、该区块的 QC、已宣布的强制指令或更新的状态映射，都不够。Ready 启动本地进入时钟；远端提案时间、私下组装 QC 的时间及 L1 安装时间都不启动它。$i'$ 的成员在进入前处理边界验证包。会话 $r+1$ 的其他消息等待进入后处理。

\subsection{待处理请求}\label{sec:requests}
每个节点维护可准入客户端请求的待处理池。只有当请求进入已封定的保留前缀后，它才离开待处理池。仅有区域提交不会移除它。被丢弃输出中的请求仍留在池中。接任区域的领导者像处理其他待处理请求一样包含它（\cref{sec:progress}）。

\section{跨区域协议}\label{sec:xprotocols}
每个节点运行\cref{alg:external}。接任区域的成员还运行\cref{alg:bootstrap}。相对 v3.3，这两个协议都是新增的。\Cref{alg:external} 全部标红。\Cref{alg:bootstrap} 复用 v3.3 恢复协议的第 3--5 步，红色标记相对这些步骤的差异。

\begin{protocol}{alg:external}{监视 L1 和数据流，并封定每个边界}
\color{rrRed}
每个节点 $u$ 运行以下步骤。其当前外部会话为 $r$。
\begin{enumerate}[label=\textbf{\arabic*)},itemsep=3pt]
\item \textbf{Watch（监视）。}持续执行以下操作。
\begin{itemize}
\item 从骨干网络接收会话 $r$ 的已提交区块及区域回执。获取缺失区间。
\item 读取已经终局化的 L1，并遵循控制顺序。
\item 验证并保存交接可能需要的区块、回执及 L1 证明。潜在接任区域还缓存提案及 QC，为交接准备。
\item 会话 $r$ 处于 Open 时，将其持有有效回执、且高度不超过 $H_r$ 的每个 DATA 区块视为暂定输出。
\end{itemize}
按控制顺序处理第 2、3 步。在同一位置，第 2 步优先。

\item \textbf{Force（强制）。}目标为 $r$ 的强制指令生效，即 L1 区块 $\lambda_a$ 终局化时，运行本步骤。
\begin{itemize}
\item 若会话 $r$ 已封定，只应用状态映射，然后停止。
\item 设置 $\End(r)=h_*$，将会话 $r$ 标记为 Truncated。只执行一次。
\item 丢弃高于 $h_*$ 的暂定输出，不论它携带什么 L1 引用。其请求仍待处理。
\item 获取高度 $h_*$ 处的 DATA 区块、回执及祖先链。根据保留前缀重建暂定状态。
\item 计算 $\lambda_s$ 及接任区域（\cref{sec:witness}）。在数据和接任区域都存在之前，保留截断点并等待。
\item 随后用 FORCE 见证封定会话 $r$（第 4 步）。
\end{itemize}

\item \textbf{正常完成。}会话 $r$ 处于 Open 时，有效回执 $\FC_r(T_r)$ 进入控制顺序，则运行本步骤。
\begin{itemize}
\item 检查 $T_r$ 的祖先链中，$A_r$ 之后恰有 $\TB_r$ 个 DATA 区块。
\item 计算 $\lambda_s\ge\lambda(T_r)$ 及接任区域。若没有其他 live 区域，保留 $T_r$ 为候选锚点并等待。若期间有强制指令生效，则由第 2 步接管。
\item 随后用 REGULAR 见证封定会话 $r$（第 4 步）。
\end{itemize}

\item \textbf{Seal（封定）。}由第 2 或第 3 步调用本步骤，或在收到会话 $r$ 的有效已封定见证 $\mathcal W$ 时运行。
\begin{itemize}
\item 固定 $\End(r)$、锚点 $A_{r+1}$、接任区域 $i_{r+1}$ 和 $\lambda_s$。它们此后永不改变。
\item 将截至 $A_{r+1}$ 的保留前缀追加到外部链。从待处理池移除其中的请求。
\item 将当前外部会话推进到 $r+1$。
\item 若 $u$ 是 $i_{r+1}$ 的成员，则启动\cref{alg:bootstrap}。
\end{itemize}
\end{enumerate}
\end{protocol}

每当新的区块、回执、授权或 L1 更新到达时，第 2--4 步重新运行，包括区域引擎被冻结时。这些步骤是幂等的。每个节点都会得到相同边界，因为 L1、控制顺序及选择规则是共同的。接任区域在进入之前应用所有在 $\lambda_s$ 处应生效的 L1 控制。会话 $0$ 通过其合成见证从第 4 步开始。

\begin{protocol}{alg:bootstrap}{引导接任区域}
\RR{区域 $i_{r+1}$ 的}副本 $v$ \RR{针对会话 $r$ 的已封定见证 $\mathcal W$}运行以下步骤。这里 $\Delta$ 和 $F$ 指区域 $i_{r+1}$ 的参数。
\begin{enumerate}[label=\textbf{\arabic*)},itemsep=3pt]
\item \RR{\textbf{Ready。}发生 $\tagg{Ready}_v(\mathcal W)$ 时（\cref{def:ready}），向 $i_{r+1}$ 的所有成员转发 $\mathcal W$ 及其完整验证包。接收者对该验证包的验证完成，构成它自己的 Ready 事件。它立即运行后续步骤，不等待骨干网络或领导者消息。}

\item \textbf{Status 与进入。}\RR{安装会话 $r+1$，使用锚点 $A_{r+1}$、授权 $\tagg{AUTH}_{r+1,0}$、配额 $\TB_{r+1}$ 和起始 L1 高度 $\lambda_s$。}然后：
\begin{itemize}
\item 将进入锁 $Q_v$ 设为\RR{基础依据 $Q^0_{r+1}$}。
\item 以未冻结状态进入\RR{$(r+1,0,0)$}，启动其进度检查。\RR{只设置一次 $e_v(r+1,0,0)$。}
\item 向 $L(r+1,0,0)$ 发送 $\langle\tagg{STATUS},\RR{r+1,0,0},Q_v,E\rangle_v$ 及其验证包。这里\RR{$E=\mathcal W$} 是进入证明。
\item 处理为\RR{会话 $r+1$ 的代次 $0$} 缓存的证明。
\end{itemize}
\RR{重复见证从不重置进入时间、配额或锁。}

\item \textbf{New-view。}仅在 $L(r+1,0,0)$ 运行本步骤。自其自己的进入时间起等待 $2\Delta$。若\RR{$(r+1,0,0)$} 仍是当前视图且未被冻结，则广播
$\langle\tagg{NEWVIEW},\RR{r+1,0,0,A_{r+1},Q^0_{r+1},\mathcal W}\rangle_{L(r+1,0,0)}$
及其验证包。\RR{起始区块必须恰好是 $A_{r+1}$。}

\item \textbf{Anchor vote。}为\RR{会话 $r+1$ 的代次 $0$，以 $E=\mathcal W$} 运行\cref{alg:recovery} 第 5 步。领导者获得 ANCHOR QC 后，从 $j=0$ 开始运行\cref{alg:steady} 的 Propose\RR{，DATA 高度为 $h(A_{r+1})+1$，物理位置为 $1$。会话 $r+1$ 中后续故障使用\cref{alg:recovery}。}
\end{enumerate}
\end{protocol}

\paragraph{引导时序。}
令 $e$ 为 $i_{r+1}$ 中第一个正确成员达到 Ready 的时间。第 1 步的转发使每个正确成员在 $e+\Delta$ 之前进入。假设初始领导者正确，且没有退出。领导者在 $e+3\Delta$ 之前发送 NEWVIEW；ANCHOR QC 在 $e+5\Delta$ 之前到达它；第一个 DATA 提案在 $e+6\Delta$ 之前到达每个正确成员。这不晚于每个成员的首次进度检查，其时刻为 $e_v(r+1,0,0)+6\Delta\ge e+6\Delta$。Ready 之前的骨干网络延迟没有固定上界。

\section{示例}\label{sec:examples}
\Cref{fig:endings} 使用贯穿全文的示例：$h(A_r)=100$、$\TB_r=10$、$K=4$，因此 $H_r=110$。

\paragraph{经过 DRAIN 后正常结束。}
代次 $0$ 有三个任期：DATA 101--104、105--108 和 109--110。DATA 110 的 QC 不启动下一任期。若 DATA 110 的计时器到期，则它提交，其回执使会话完成。反过来，假设退出先取消了这些计时器。代次 $1$ 随后从已认证的 DATA 110 开始，因此是收尾代次。它唯一的新区块是一个 DRAIN。DRAIN 的计时器将 DATA 110 作为祖先提交。副本为 DATA 110 发送 FINAL，不为 DRAIN 发送。若再次退出，则需要再增加一个代次和一个 DRAIN。无论哪种情况，会话 $r+1$ 都从 DATA 111 开始，其父区块为 DATA 110。

\paragraph{强制结束。}
委员会发布目标为会话 $r$ 的强制指令，其中 $h_*=106$、$\lambda_a=1200$。旧区域仍可能在激活之前或之后提交 DATA 107 和 108。一旦 L1 区块 1200 终局化，每个节点仅将会话 $r$ 保留至 DATA 106。它丢弃高于 106 的暂定输出，其中的请求仍待处理。验证 DATA 106 及其回执之后，接任区域开启 $(r+1,0,0)$。它提出自己的 DATA 107，父区块为 DATA 106。旧区域日志不改变。其 107 和 108 的回执永远不能将这些区块带回外部链。

\begin{figure}[htbp]\centering
% Two ways a session ends (Section 7), with h(A_r)=100, TB_r=10, K=4.
% (a) 经过一个 DRAIN 后正常结束.  (b) Forced ending at h_*=106.
% Teal: 外部链. Gray: 仅在区域日志中. Blue: session r+1.
% Arrows point from parent to child.
\begin{tikzpicture}[x=1cm,y=1cm,font=\fontsize{9}{10.5}\selectfont,
  blk/.style={draw,rounded corners=1.5pt,minimum height=.62cm,
    minimum width=.95cm,inner sep=2pt},
  ext/.style={blk,draw=rrTeal,fill=rrTeal!12},
  reg/.style={blk,draw=black!45,fill=black!6,text=black!65},
  nxt/.style={blk,draw=rrBlue,fill=rrBlue!10},
  link/.style={-{Stealth[length=1.4mm]},semithick,draw=black!70},
  lane/.style={anchor=east,font=\fontsize{8}{9}\selectfont,text=black!70},
  tag/.style={font=\fontsize{8}{9}\selectfont,inner sep=1pt},
  ttl/.style={anchor=west,font=\bfseries\fontsize{9}{10.5}\selectfont},
  note/.style={anchor=north west,align=flush left,text width=5.7cm,
    inner sep=0pt,font=\fontsize{8.5}{10}\selectfont},
  brace/.style={decorate,decoration={brace,amplitude=3pt},draw=black!55}]
  % ---------------- (a) regular ending ----------------
  \node[ttl] at (-1.55,1.3) {(a) 经过一个 DRAIN 后正常结束};
  \node[lane] at (.35,0) {会话 $r$};
  \node[lane] at (.35,-1.2) {会话 $r{+}1$};
  \node[ext] (a100) at (.95,0) {100};
  \node[ext,minimum width=1.5cm] (a101) at (2.55,0) {101--104};
  \node[ext,minimum width=1.5cm] (a105) at (4.35,0) {105--108};
  \node[ext] (a109) at (5.8,0) {109};
  \node[ext] (a110) at (7.0,0) {110};
  \node[reg,dashed] (aD) at (8.35,0) {DRAIN};
  \node[nxt] (a111) at (8.35,-1.2) {111};
  \draw[link] (a100)--(a101); \draw[link] (a101)--(a105);
  \draw[link] (a105)--(a109); \draw[link] (a109)--(a110);
  \draw[link] (a110)--(aD);
  \draw[link] (a110.south)--(a111.west);
  \node[tag,below=2pt of a100] {$A_r$};
  \node[tag,anchor=north east] at ([xshift=-1pt,yshift=-2pt]a110.south west) {$T_r$};
  \node[tag,above=3pt of a101] {任期 0};
  \node[tag,above=3pt of a105] {任期 1};
  \draw[brace] ([yshift=3pt]a109.north west)--([yshift=3pt]a110.north east)
    node[tag,midway,above=4pt] {任期 2};
  \node[tag,above=3pt of aD] {代次 1};
  \node[note] at (9.4,.5) {退出在 110 提交之前取消其计时器。代次 1 从 110 开始，因此是收尾代次。其唯一 DRAIN 的计时器提交 110，且仅为 110 发送 FINAL。会话 $r{+}1$ 从 111 开始，父区块为 110。};
  % ---------------- (b) forced ending ----------------
  \def\yb{-3.65}\def\yn{-4.85}
  \node[ttl] at (-1.55,\yb+1.3) {(b) 在 $h_*=106$ 处强制结束};
  \node[lane] at (.35,\yb) {会话 $r$};
  \node[lane] at (.35,\yn) {会话 $r{+}1$};
  \node[ext] (b100) at (.95,\yb) {100};
  \node[ext,minimum width=1.5cm] (b101) at (2.55,\yb) {101--104};
  \node[ext] (b105) at (4.0,\yb) {105};
  \node[ext] (b106) at (5.2,\yb) {106};
  \node[reg] (b107) at (6.65,\yb) {107};
  \node[reg] (b108) at (7.85,\yb) {108};
  \node[nxt] (n107) at (6.65,\yn) {107};
  \draw[link] (b100)--(b101); \draw[link] (b101)--(b105);
  \draw[link] (b105)--(b106); \draw[link] (b106)--(b107);
  \draw[link] (b107)--(b108);
  \draw[link] (b106.south)--(n107.west);
  \draw[rrRed,thick,densely dashed] (5.92,\yb+.55)--(5.92,\yb-.5);
  \node[tag,text=rrRed,above] at (5.92,\yb+.55) {截断点};
  \node[tag,below=2pt of b100] {$A_r$};
  \node[tag,anchor=north east] at ([xshift=-1pt,yshift=-2pt]b106.south west) {$h_*$};
  \node[note] at (9.4,\yb+.5) {强制指令在 L1 高度 1200 激活。外部链仅将会话 $r$ 保留至 106。旧区域仍可能提交 107、108；外部链丢弃它们。会话 $r{+}1$ 提出自己的 DATA 107，父区块为 106。};
  % ---------------- legend ----------------
  \def\yl{-5.95}
  \node[ext,minimum width=.5cm,minimum height=.34cm] at (-.9,\yl) {};
  \node[anchor=west,tag] at (-.55,\yl) {外部链};
  \node[reg,minimum width=.5cm,minimum height=.34cm] at (2.3,\yl) {};
  \node[anchor=west,tag] at (2.65,\yl) {仅在区域日志中};
  \node[nxt,minimum width=.5cm,minimum height=.34cm] at (5.7,\yl) {};
  \node[anchor=west,tag] at (6.05,\yl) {会话 $r{+}1$};
\end{tikzpicture}

\caption{会话结束的两种方式，其中 $h(A_r)=100$、$\TB_r=10$、$K=4$。箭头从父区块指向子区块。(a) 退出使 DATA 110 已获认证但尚未提交；一个 DRAIN 使其提交。(b) 强制指令将外部链截断于 106；旧区域的 107、108 只保留在其日志中。}\label{fig:endings}
\end{figure}

\begingroup\small
\begin{thebibliography}{9}
\bibitem{rr33}
Sync HotStuff RR v3.3：提案计时器与末位 QC 交接（Proposal Timers and Terminal-QC Handoff）。规范草案，2026 年 10 月 12 日（\texttt{sync\_hotstuff\_rr\_v3\_3.tex}）。
\bibitem{sync}
I. Abraham, D. Malkhi, K. Nayak, L. Ren, and M. Yin.
\emph{Sync HotStuff: Simple and Practical Synchronous State Machine Replication}
（Sync HotStuff：简单实用的同步状态机复制）。
IEEE Symposium on Security and Privacy, 2020, pp. 106--118.
\href{https://conferences.computer.org/sp/pdfs/sp/2020/349700a106.pdf}{已发表论文}。
\end{thebibliography}
\endgroup
\end{document}
